% This is the first rough structure for my conference paper.
% I used the official IEEE Conference Template on Overleaf as a guide
% when deciding on the initial layout and structure of this file.
% NB ** Final paper must use the IEEE conference format and be exactly 2 pages.




\documentclass[conference]{IEEEtran}




\usepackage{cite}




\title{Closing the Coverage Gap with AI: AI-Assisted Code Coverage Management for Embedded C Software}




\author{
\IEEEauthorblockN{Dara Heaphy}
\IEEEauthorblockA{
University of Limerick\\
Limerick, Ireland
}
}




\begin{document}




\maketitle




\section{Introduction}




% Going to start broad here.
% Briefly introduce embedded software testing, code coverage,
% and why improving coverage feedback could be useful.



Embedded software is used in systems where reliable behaviour is often
important, which makes testing a major part of the development process.
Code coverage is commonly used to show which parts of a program have or
have not been exercised by a test suite, but a coverage report on its own
does not explain why code remains uncovered or what should be done next.
This can be particularly difficult in embedded C, where hardware
dependencies, peripheral interactions, timing behaviour, and the need for
mocks or stubs can make certain paths harder to test \cite{garousi2018}.
With large language models now being explored for software testing and test
generation, there may be an opportunity to make this part of the coverage
process more useful by using AI to help explain uncovered code and suggest
possible ways of exercising it.




% Mini literature review.
% I want to group the papers by idea rather than go through them one by one.
% Main areas:
% - LLM-based test generation
% - coverage-guided feedback
% - static/program analysis
% - embedded C testing, mocks and stubs



Recent work such as CoverUp, Panta, and SymPrompt has shown that LLM-based
test generation can be improved by using coverage feedback, program analysis,
and path-specific context \cite{pizzorno2025,gu2026,ryan2024}. However, these
approaches are mainly evaluated outside embedded C. Existing embedded testing
work such as SmartUnit and CTGEN already deals with automated testing, C code,
and the use of stubs \cite{zhang2018,mangels2012}, but does not focus on using
AI to explain unresolved coverage gaps and suggest practical tests, mocks, or
stubs.


A recurring theme across the more recent work is that LLM-based test
generation becomes more useful when the model is given focused information
about the code that remains difficult to cover. HITS uses method slicing to
reduce complex testing problems into smaller targets, while static-analysis
guided approaches provide selected dependency and program information rather
than an entire codebase \cite{wang2024hits,chowdhury2024}. CodaMOSA similarly
introduces an LLM when conventional automated test generation reaches a
coverage plateau, and ChatUniTest combines generation with validation and
repair of unsuccessful tests \cite{lemieux2023,chen2024}. Together, these
approaches suggest that an embedded-C assistant should analyse a specific
uncovered region using relevant branch conditions, dependencies, existing
tests, and hardware-abstraction information instead of simply being asked to
generate additional tests from the source code alone.




% Research gap.
% The main idea at the moment is that most LLM testing work is not focused
% on embedded C, while embedded testing research does not really look at
% AI-assisted explanation of coverage gaps and suggested tests/mocks/stubs.


This leaves a more specific gap between current LLM-based testing research
and established embedded-software testing techniques. Embedded testing
research already recognises the importance of hardware dependencies,
abstraction layers, stubs, and other constraints that affect whether code can
be exercised realistically \cite{garousi2018,zhang2018,mangels2012}, but these
issues have not been the main focus of the LLM approaches reviewed here. This
study therefore asks to what extent AI-assisted coverage analysis can improve
the identification, explanation, and reduction of coverage gaps in embedded C
software. Baseline line, branch, and function coverage will first be collected,
after which selected uncovered regions will be provided to an AI model together
with relevant source and environment context. Suggested tests, mocks, or stubs
will then be manually reviewed, compiled, executed, and compared against the
baseline, with technical correctness and usefulness considered alongside any
increase in coverage \cite{inozemtseva2014}. The intended contribution of this
study is an evaluation of whether an AI-assisted feedback step can make
embedded-C coverage reports more actionable, including evidence about the
technical correctness and practical usefulness of suggested tests, mocks, and
stubs.




% My contribution.
% Explain the main research question and briefly describe the approach:
% baseline coverage -> AI analysis -> review suggestions -> run selected
% changes -> compare coverage and assess whether the suggestions were useful.




% Finish with a short paper roadmap once the rest of the paper structure
% is properly decided.


The remainder of the proposed paper is structured as follows. Section II
describes the AI-assisted coverage workflow and evaluation methodology,
Section III presents the results, and Section IV discusses the findings,
limitations, and conclusions.




\bibliographystyle{IEEEtran}
\bibliography{references}




\end{document}
